\section*{Capítulo \ref{ch:b1:chemical-biogenic-sediments} --- Sedimentos carbonatados, silíceos y evaporíticos}

\begin{solution}{exo:b1:chemical-biogenic-sediments:1}
Cloruro $0.5503\times35.165 = \qty{19.35}{g}$; sodio $0.3066\times35.165 =
\qty{10.78}{g}$; sulfato \qty{2.711}{g}; magnesio \qty{1.284}{g}. El sodio y el cloruro juntos suman $0.5503 + 0.3066 = 0.857$, es decir, el \qty{85.7}{\%} de
la sal.
\end{solution}

\begin{solution}{exo:b1:chemical-biogenic-sediments:2}
La fracción que queda es $1/\mathrm{CF}$: carbonato de calcio $1/1.8 = 0.556$ queda, \qty{44.4}{\%} evaporado; yeso $0.263$ queda, \qty{73.7}{\%} evaporado;
halita $0.094$ queda, \qty{90.6}{\%} evaporada; sulfato de magnesio $0.014$ queda, \qty{98.6}{\%} evaporado.
\end{solution}

\begin{solution}{exo:b1:chemical-biogenic-sediments:3}
(a)~$\rho_s(1-\phi) = G/a = 8.4/0.0042 = \qty{2000}{kg.m^{-3}}$. (b)~Con aragonita, $1 - 2000/2930 = 0.32$, una porosidad del \qty{32}{\%}; con calcita,
$1 - 2000/2710 = 0.26$. (c)~$a = 3.7/2000 = \qty{1.85e-3}{m}$, \qty{1.85}{mm} al año.
\end{solution}

\begin{solution}{exo:b1:chemical-biogenic-sediments:4}
Hasta \qty{4000}{m} el contenido baja solo lentamente (de 88 a 80); por debajo de \qty{4000}{m} baja bruscamente: la \omterm{def:b1:chemical-biogenic-sediments:ccd}{lisoclina} yace a unos \qty{4000}{m}. El contenido
baja por debajo del \qty{10}{\%} entre \qty{4700}{m} (12) y \qty{4800}{m} (6): $4700 + 100\times(12-10)/(12-6) = \qty{4733}{m}$, una CCD a unos \qty{4730}{m}, unos
\qty{700}{m} por debajo de la \omterm{def:b1:chemical-biogenic-sediments:ccd}{lisoclina}.
\end{solution}

\begin{solution}{exo:b1:chemical-biogenic-sediments:5}
(a)~\ce{Ca^2+ + 2 HCO3- -> CaCO3 + CO2 + H2O}: una mol de \ce{CO2} por mol de \ce{CaCO3}, así que $44.01/100.09 = \qty{0.440}{kg}$ de \ce{CO2} por kilogramo de
\ce{CaCO3}. (b)~$3.7\times0.440 = \qty{1.63}{kg}$ de \ce{CO2} por metro cuadrado y por año. (c)~De los dos iones hidrogenocarbonato usados, solo un carbono entra en el
\omterm{def:b1:rock-forming-minerals:mineral}{mineral}; el otro sale como dióxido de carbono. La calcificación saca alcalinidad del agua y reduce su capacidad de retener carbono disuelto; cuánto del \ce{CO2} llega al aire
depende del poder tampón del agua de mar, asunto de los equilibrios carbonatados de los libros de química.
\end{solution}

\begin{solution}{exo:b1:chemical-biogenic-sediments:6}
(a)~$t_c = (2000/350)^2 = \qty{32.7}{Myr}$; a \qty{25}{km} por millón de años, $25\times32.7 = \qty{816}{km}$ del eje. (b)~$t_c = (2600/350)^2 =
\qty{55.2}{Myr}$, \qty{1380}{km}. (c)~Un fondo de \qty{60}{Myr} ha recibido arcilla roja durante $60 - 32.7 = \qty{27.3}{Myr}$ en el primer caso, y solo durante
$60 - 55.2 = \qty{4.8}{Myr}$ en el segundo.
\end{solution}

\begin{solution}{exo:b1:chemical-biogenic-sediments:7}
(a)~Sí: $a = 4.2 > r = 3$, y el \omterm{def:b1:chemical-biogenic-sediments:reef}{arrecife} gana \qty{1.2}{mm} al año sobre el mar hasta que su cima llega a la superficie. (b)~$D$ crece $10 - 4.2 = \qty{5.8}{mm}$
al año; de 2 a \qty{50}{m} tarda $48/0.0058 = \num{8300}$ años. (c)~Porque su tasa de crecimiento baja a medida que el agua se hace más profunda: llega menos luz a las
algas que viven en los corales y los alimentan, así que $a$ disminuye y el \omterm{def:b1:chemical-biogenic-sediments:reef}{arrecife} se queda atrás cada vez más deprisa.
\end{solution}

\begin{solution}{exo:b1:chemical-biogenic-sediments:8}
(a)~La reacción es \ce{CaSO4.2H2O -> CaSO4 + 2 H2O}. (b)~Los volúmenes molares son $172.17/2.30 =
\qty{74.9}{cm^3}$ y $136.14/2.97 = \qty{45.8}{cm^3}$, así que el sólido pierde $1 - 45.8/74.9 = \qty{38.8}{\%}$ de su volumen. (c)~$2\times18.02/172.17 = 0.209$:
\qty{20.9}{\%} de la masa sale como agua. Liberada a profundidad, esta agua eleva la presión en los poros de las rocas circundantes; bajo un sello puede debilitar las rocas y ayudar
a que las fallas deslicen, y puede disolver y transportar otras sales.
\end{solution}

\begin{solution}{exo:b1:chemical-biogenic-sediments:9}
(a)~$9.2/255 = \qty{3.6}{\%}$ conservado, \qty{96.4}{\%} disuelto. (b)~Donde la producción es muy alta, un pequeño porcentaje de ella sigue siendo un flujo grande, y donde
poco barro o carbonato lo diluyen: bajo las aguas de surgencia del océano Austral, en el Pacífico norte y en el Pacífico ecuatorial, y por debajo de la CCD. (c)~Las frústulas finas y
delicadas se disuelven primero y las robustas sobreviven: el sedimento conserva una selección de las especies más resistentes, no la comunidad que vivió en el agua.
\end{solution}

\begin{solution}{exo:b1:chemical-biogenic-sediments:10}
(a)~$EA = 1.5\times\qty{1e10}{m^2} = \qty{1.5e10}{m^3}$ al año. (b)~Con
$Q_{\mathrm{in}} = \qty{1.8e10}{m^3}$, $Q_{\mathrm{out}} = \qty{0.3e10}{m^3}$ y
$k = 6$: por encima del umbral del yeso (3.8) pero no del de la halita (10.6), así que precipitan el carbonato de calcio y el yeso. Con $Q_{\mathrm{in}} =
\qty{1.575e10}{m^3}$, $Q_{\mathrm{out}} = \qty{0.075e10}{m^3}$ y $k = 21$, para un ion que permanece en disolución como el magnesio: la salmuera pasa el umbral de la halita, así que
precipita también la halita, pero no las sales de magnesio (70). (c)~La halita requiere $Q_{\mathrm{in}}/(Q_{\mathrm{in}} - EA) \ge 10.6$, es decir
$Q_{\mathrm{in}} \le 10.6\,EA/9.6 = 1.104\,EA = \qty{1.66e10}{m^3}$ al año: la entrada puede exceder la evaporación en a lo sumo una décima parte.
\end{solution}

\begin{solution}{exo:b1:chemical-biogenic-sediments:11}
(a)~El agua de mar contiene \qty{67.3}{mg} de bromuro por kilogramo; a 10.6~veces, la salmuera contiene $67.3\times10.6 = \qty{713}{mg.kg^{-1}}$ y la primera halita
$0.032\times713 = \qty{22.8}{mg.kg^{-1}}$. (b)~A 30~veces, $67.3\times30 =
\qty{2019}{mg.kg^{-1}}$ en la salmuera y \qty{64.6}{mg.kg^{-1}} en la halita. (c)~La halita toma solo 3.2~centésimas de la concentración de la salmuera, así que casi todo el bromuro
permanece en disolución y sigue concentrándose a medida que el agua se evapora. (d)~La subida registra la concentración de la salmuera; la caída brusca registra su dilución, por una nueva
entrada de agua de mar en la cuenca, o de agua dulce, antes de que la halita siguiera creciendo.
\end{solution}

\begin{solution}{exo:b1:chemical-biogenic-sediments:12}
(a)~$t_c = \qty{32.65}{Myr}$, así que $0.02\times\num{32650} = \qty{653}{m}$ de fango (\num{32650} intervalos de mil años a \qty{2}{cm}) y
$0.001\times\num{37350} = \qty{37}{m}$ de arcilla roja. (b)~$t_c = (2500/350)^2 =
\qty{51.0}{Myr}$: \qty{1020}{m} de fango y \qty{19}{m} de arcilla roja. Una CCD más profunda da un fango mucho más grueso, y la arcilla roja del fondo atlántico profundo es más
delgada que la del Pacífico de la misma edad. (c)~La CCD pudo ser más profunda durante un tiempo, de modo que el emplazamiento subió de nuevo por encima de ella; o la placa llevó el emplazamiento
bajo un cinturón de alta productividad, como el ecuador, donde llovía más carbonato; o la capa la trajo un deslizamiento o una \omterm{def:b1:clastic-sediments:turbidite}{corriente de turbidez} desde una parte más alta del fondo marino.
\end{solution}

\begin{solution}{pb:b1:chemical-biogenic-sediments:1}
\textbf{1.} $1030\times1000 = \qty{1.03e6}{kg}$ de agua de mar, que contienen
$1.03\times10^6\times0.035165 = \qty{36200}{kg}$ de sal.

\textbf{2.} Yeso $1.03\times10^6\times1.581\times10^{-3} = \qty{1628}{kg}$;
halita $1.03\times10^6\times0.02741 = \qty{28230}{kg}$.

\textbf{3.} Carbonato de calcio $113.3/2710 = \qty{0.042}{m}$; yeso $1628/2300 =
\qty{0.71}{m}$; halita $28230/2170 = \qty{13.0}{m}$.

\textbf{4.} La anhidrita tiene la masa $1628\times136.14/172.17 =
\qty{1288}{kg}$ y el espesor $1288/2970 = \qty{0.43}{m}$.

\textbf{5.} Secas, las sales suman $0.110 + 1.250 + 27.41 + 6.34 = \qty{35.11}{g}$ por kilogramo, de las cuales la halita es $27.41/35.11 = \qty{78}{\%}$: el sodio y el cloruro
forman la mayor parte de la sal disuelta, y la halita se lleva casi todos ellos.

\textbf{6.} De abajo arriba: carbonato de calcio, yeso, halita, las sales amargas. Cada \omterm{def:b1:rock-forming-minerals:mineral}{mineral} empieza a precipitar cuando la salmuera alcanza su saturación, el menos soluble primero, y
se deposita sobre lo que se depositó antes.

\textbf{7.} La profundidad de salmuera es $1000/\mathrm{CF}$: \qty{556}{m} para el carbonato de calcio, \qty{263}{m} para el yeso, \qty{94}{m} para la halita.

\textbf{8.} $1 - 1/10.6 = \qty{90.6}{\%}$.

\textbf{9.} $1000/70 = \qty{14}{m}$ para el sulfato de magnesio, $1000/90 = \qty{11}{m}$ para las sales de potasio.

\textbf{10.} Unos \qty{1.10}{g.cm^{-3}} cuando se forma el carbonato de calcio, 1.13 con el yeso, 1.22 con la halita y más de 1.30 con las sales amargas, frente a 1.03 del agua de mar. El agua
de mar que entra flota sobre la salmuera densa, que se queda en el fondo; en una cuenca con umbral la salmuera puede volver al océano por debajo de la entrada.

\textbf{11.} Solo tiene \qty{4}{cm} de espesor, frente a \qty{13}{m} de halita; y se forma en los márgenes, donde el agua que entra se concentra primero, más que en el centro, así que se mezcla con la base del yeso o queda sustituida por él.

\textbf{12.} La halita toma solo una pequeña parte del bromuro de la salmuera (ejercicio~\ref{exo:b1:chemical-biogenic-sediments:11}), así que la salmuera se enriquece en bromuro a medida que crece la halita, y cada nueva capa de halita toma más.

\textbf{13.} $1500/13.01 = 115.3$: 116 llenados.

\textbf{14.} $116\times0.434 = \qty{50}{m}$ de anhidrita, no como un estrato, sino como 116 capas, cada una algo menor de medio metro de espesor, en la base de la halita de su llenado: una sucesión cíclica de anhidrita y halita.

\textbf{15.} Por llenado, $1.03\times10^6\times6.34\times10^{-3}/2000 =
\qty{3.3}{m}$, así que \qty{379}{m} para 116 llenados. Son las sales más solubles: la siguiente entrada de agua de mar las redisuelve, y una cuenca cuya salmuera sale antes del final nunca llega a ellas.

\textbf{16.} En cada rellenado la salmuera se diluye: el contenido de bromuro cae de golpe en la base de la nueva halita y vuelve a subir hacia arriba, un perfil en dientes de sierra con un diente por llenado.

\textbf{17.} $2000/13.01 = 154$ llenados.

\textbf{18.} $5000\times365.25\times86400 = \qty{1.578e11}{m^3}$ al año, repartidos sobre $\qty{1.0e11}{m^2}$: $E = \qty{1.58}{m}$ al año.

\textbf{19.} $1.578\times10^{11}\times1030\times0.035165 = \qty{5.72e12}{kg}$ de sal al año.

\textbf{20.} $1.578\times10^{11}\times1030\times0.02741/2170 =
\qty{2.05e9}{m^3}$ de halita al año, es decir \qty{0.0205}{m}, unos \qty{2}{cm} al año sobre el fondo.

\textbf{21.} $1500/0.0205 = \num{73000}$ años.

\textbf{22.} El nivel baja \qty{1.58}{m} al año; se van \qty{1000}{m} de agua en $1000/1.578 = 634$ años.

\textbf{23.} La halita aparece cuando se ha ido el \qty{90.6}{\%} del agua, tras $905.7/1.578 = 574$ años.

\textbf{24.} Los \qty{1.5}{km} de halita tomaron la sal de 116 llenados de la cuenca (115.3, que hay que redondear hacia arriba), ya sea que llegaran uno tras otro al secarse y rellenarse la cuenca, o de forma continua
por el umbral, como en la pregunta~21, a razón de un volumen de agua de mar cada $1000/1.578 = 634$ años.
\end{solution}
